\chapter{Conclusioni e Sviluppi Futuri}
\label{chap:conclusioni}

\section{Obiettivi raggiunti}
Il prototipo \textit{EscapeManager} ha centrato gli obiettivi iniziali:
\begin{itemize}
    \item \textbf{Architettura solida:} l'applicazione è strutturata secondo il paradigma BCE (Boundary-Control-Entity) con una rigida separazione delle responsabilità tra i livelli di presentazione, controllo, dominio e persistenza.
    \item \textbf{Design Pattern GoF:} sono stati applicati e documentati 5 pattern fondamentali: \textbf{State} (ciclo di vita della Stanza), \textbf{Strategy} (politiche di prezzo), \textbf{Observer} (lista d'attesa), \textbf{Builder} (costruzione fluente delle Prenotazioni) e \textbf{Singleton} (gestione centralizzata della connessione DB).
    \item \textbf{Pattern architetturali:} DAO, Abstract Factory (DAO Factory) e Dependency Injection sono utilizzati per disaccoppiare la logica di business dalla tecnologia di persistenza.
    \item \textbf{Persistenza relazionale:} il database PostgreSQL, progettato con Single Table Inheritance e junction table, supporta integralmente il Domain Model.
    \item \textbf{Testing multi-livello:} la suite di collaudo comprende test unitari (White-Box), test con mock (Mockito), test di integrazione (DAO) e test funzionali (Black-Box) tracciati sui casi d'uso, per un totale di 44 test cases.
\end{itemize}

\section{Sviluppi Futuri}

Essendo suddiviso in moduli indipendenti, il sistema lascia aperte molte possibilità per eventuali estensioni. Di seguito alcune idee per tappe future (funzionalità non implementate in questa prima versione):

\subsection{Interfaccia Grafica Completa (JavaFX/Web)}
I mockup attuali possono essere trasformati in una GUI funzionale collegata ai Controller esistenti, senza modificare il Domain Model. In alternativa, una migrazione a frontend web (React/Angular) con backend Spring Boot richiederebbe solo la trasformazione dei Controller in \texttt{@RestController}.

\subsection{Autenticazione e Autorizzazione}
Il campo \texttt{password\_hash} è già predisposto nel modello dati. L'implementazione prevede:
\begin{itemize}
    \item Hashing sicuro delle password con BCrypt o Argon2
    \item Login con generazione di token JWT (JSON Web Token)
    \item Controllo accessi basato su ruolo (\texttt{RBAC}) con annotazioni \texttt{@Secured}
    \item Session management con timeout automatico
\end{itemize}

\subsection{Notifiche Push per Lista d'Attesa}
L'infrastruttura Observer è implementata a livello di dominio. L'estensione futura prevede:
\begin{itemize}
    \item Invio email automatico tramite JavaMail API quando uno slot si libera
    \item Notifiche push mobile tramite Firebase Cloud Messaging (FCM)
    \item SMS tramite Twilio API per notifiche urgenti
    \item Dashboard in tempo reale con WebSocket per aggiornamenti live
\end{itemize}

\subsection{Integrazione con Sistemi di Pagamento}
Integrazione con gateway come Stripe o PayPal per gestire pagamenti online:
\begin{itemize}
    \item Persistenza dello stato transazionale (PENDING, COMPLETED, FAILED)
    \item Rollback automatico della prenotazione in caso di fallimento pagamento
    \item Refund automatico in caso di cancellazione entro X ore
    \item Gestione 3D Secure per transazioni sicure
\end{itemize}

\subsection{Integrazione Google Calendar e Outlook}
Sincronizzazione bidirezionale delle prenotazioni con calendar esterni:
\begin{itemize}
    \item API OAuth2 per accesso sicuro ai calendar
    \item Creazione automatica eventi in Google Calendar alla conferma prenotazione
    \item Reminder automatici via email 24h prima della sessione
    \item iCalendar export per compatibilità con tutti i client
\end{itemize}

\subsection{Sistema di Raccomandazioni con Machine Learning}
Analisi dei pattern di prenotazione per suggerimenti personalizzati:
\begin{itemize}
    \item Algoritmi di Collaborative Filtering per raccomandare stanze simili
    \item Analisi predittiva della domanda per ottimizzazione prezzi dinamici
    \item Clustering dei clienti per campagne marketing mirate
    \item Sentiment analysis delle recensioni per improvement continuo
\end{itemize}

\subsection{Dashboard Analytics e Reporting}
Modulo di Business Intelligence per amministratori:
\begin{itemize}
    \item KPI in tempo reale (tasso di occupazione, revenue per stanza, NPS)
    \item Grafici interattivi con Chart.js o D3.js
    \item Export report in PDF/Excel per presentazioni executive
    \item Heatmap delle prenotazioni per identificare fasce orarie critiche
\end{itemize}

\subsection{Gestione Multi-Sede Avanzata}
L'entità \texttt{Sede} è già presente nel database. Estensioni:
\begin{itemize}
    \item Dashboard centralizzata per manager multi-filiale
    \item Trasferimento prenotazioni tra sedi (rebooking automatico)
    \item Configurazione tariffe e stati differenziati per sede
    \item Sincronizzazione inventario (es. kit di gioco) tra sedi
\end{itemize}

\subsection{Adapter Pattern per IoT e Gamification}
Integrazione di sensori fisici e meccaniche di gioco:
\begin{itemize}
    \item Timer fisico in stanza con integrazione REST API
    \item Sensori di completamento enigmi (RFID, NFC) per tracking progress
    \item Leaderboard globale con classifica performance giocatori
    \item Achievement system con badge e rewards virtuali
    \item Integrazione Arduino/Raspberry Pi per automazione stanza
\end{itemize}
